\section{稳定化与正则化策略}

\subsection{Bootstrapping 与梯度归一}

Bootstrapping 将网络输出与自身目标绑定，容易引发梯度爆炸。实践中常用梯度裁剪和归一化来控制。

\begin{importantbox}{稳定技巧速查}
\begin{itemize}
    \item 梯度裁剪到 1.0，避免更新过大导致发散。
    \item Layer normalization 和 weight decay 控制输出范围。
    \item Double Q 的动作选择与估值拆分也属于稳定策略。
\end{itemize}
\end{importantbox}

\subsection{奖励塑形与行为规范}

在高维任务中，直接使用 reward 会产生高方差。可以添加 shaping term 或 entropy penalty 平滑 learning signal。

\begin{warningbox}{奖励塑形的陷阱}
shaping term 必须是 potential-based，否则可能改变最优策略；添加 penalty 也不能让 Q 值无限增大。
\end{warningbox}

\subsection{动作约束与输出裁剪}

为了避免 Q 函数输出漂移，可采用 clipping 与 output constraint：

\begin{itemize}
    \item 将 TD target 限制在 observed reward range ± margin，防止估值爆炸。
    \item 引入 Q regularization（如 L2 shrinkage toward previous target）使更新更平滑。
    \item 对 action 采样加噪后 clamp，保证 $\arg\max_a Q(s,a)$ 不会落入未见区域。
\end{itemize}

\subsection{本章小结}

稳定化需要同时从梯度信号和行为信号两个方向着手，避免值函数爆炸与策略漂移。

